\documentclass[11pt,a4paper]{article}
\usepackage[utf8]{inputenc}
\usepackage[T1]{fontenc}
\usepackage[portuguese]{babel}
\usepackage{xcolor}
\usepackage{hyperref}
\usepackage{geometry}
\usepackage{listings}
\usepackage{tcolorbox}
\geometry{margin=2.5cm}
\definecolor{primary}{RGB}{41,128,185}
\definecolor{secondary}{RGB}{44,62,80}
\definecolor{codebg}{RGB}{240,240,240}
\lstset{backgroundcolor=\color{codebg},basicstyle=\ttfamily\small,breaklines=true,frame=single}
\title{\color{secondary}\Huge\textbf{Monitoramento de Aplicações no Kubernetes}}
\author{\color{primary}DevOps com Kubernetes -- 2026}
\date{}
\begin{document}
\maketitle


\section*{\color{primary}O que você aprenderá nesta página}
\begin{itemize}
\item Entender por que monitorar aplicações publicadas
\item Coletar métricas com Prometheus e visualizar com Grafana
\item Relacionar métricas com escalabilidade e diagnóstico
\end{itemize}

Publicar uma aplicação não termina quando o \texttt{kubectl apply} funciona. Depois que ela está em execução, precisamos responder perguntas simples e importantes: a aplicação está saudável? O consumo de CPU subiu? Os Pods estão reiniciando? O banco está respondendo? O tráfego aumentou?

Monitoramento é a prática de coletar, armazenar e visualizar sinais sobre o sistema. No Kubernetes, esses sinais vêm dos próprios recursos do cluster, dos contêineres e da aplicação.

\section*{\color{primary}Métricas do cluster}

O Metrics Server fornece métricas básicas para comandos como:

\begin{lstlisting}[language=bash]
kubectl top pods
kubectl top nodes
\end{lstlisting}

Ele é útil para uma visão rápida e também serve de base para recursos como HPA em muitos ambientes. Porém, para análises históricas e painéis completos, normalmente usamos Prometheus e Grafana.

\section*{\color{primary}Prometheus e Grafana}

Prometheus coleta métricas periodicamente a partir de endpoints HTTP. Grafana permite montar painéis para visualizar essas métricas. Em ambientes Kubernetes, uma instalação comum é feita com Helm usando o chart \texttt{kube-prometheus-stack}.

\begin{lstlisting}[language=bash]
helm repo add prometheus-community https://prometheus-community.github.io/helm-charts
helm repo update
helm install kube-prometheus-stack prometheus-community/kube-prometheus-stack -n monitoring --create-namespace
\end{lstlisting}

\section*{\color{primary}Métricas úteis para publicação}

Durante a construção e publicação de aplicações, algumas métricas são especialmente importantes: reinícios de contêineres, quantidade de Pods prontos, consumo de CPU e memória, latência, taxa de erro e disponibilidade de banco de dados. Essas métricas ajudam a decidir se uma versão pode continuar em produção ou se precisa ser revertida.

\section*{\color{primary}Boas práticas}

Configure endpoints de saúde, use readiness e liveness probes, acompanhe reinícios e mantenha dashboards simples. Um painel bom não é o que mostra tudo; é o que ajuda a perceber rapidamente quando algo importante saiu do normal.

\end{document}
